% 03 · Clustering fundamentals — what a cluster is, and what breaks
\section{CLUSTERING FUNDAMENTALS}
\label{sec:fundamentals}

\deckslides{slides 3 (the clustering problem), 18 (no free lunch) and 19--20 (validation).}

Before any algorithm, two questions deserve a straight answer: what do we
mean by ``a cluster'', and what does a distance in 16 dimensions actually
measure? The answers are not academic. Most of the failures in
\textbf{Section~\ref{sec:benchmark}} are failures of an assumption about one
of those two things, not failures of an implementation.

\subsection{Four definitions of ``cluster''}

A clustering algorithm does not discover the clusters that are in the data;
it discovers the clusters that its own definition of ``cluster'' implies. The
definitions in play in this workbook are:

\begin{itemize}
\item \textbf{Centroid.} A cluster is a set of points closer to one centre
  than to any other. This is K-means (\S\ref{sec:kmeans}), and it implies
  convex, roughly equal-sized groups.
\item \textbf{Density.} A cluster is a connected region where the local point
  density exceeds a threshold. This is DBSCAN and its descendants
  (\S\ref{sec:dbscan}--\S\ref{sec:plscan}), and it implies that clusters and
  noise are distinguishable by density alone.
\item \textbf{Connectivity.} A cluster is a connected component of a graph
  built from nearest neighbours, with edges weighted by similarity. This is
  the graph view that t-SNE, UMAP and EVoC share
  (\S\ref{sec:tsne}--\S\ref{sec:evoc}).
\item \textbf{Statistical.} A cluster is a component of a mixture model whose
  parameters are inferred by maximum likelihood (the EM algorithm,
  \citealp{Dempster:77}). The cluster boundaries are then soft: a star has a
  probability of belonging to each group.
\end{itemize}

The reference chapter behind the lecture's framing --- \citet{GarciaDias:20},
a methods chapter written for a psychiatry audience --- makes the same point
from the other direction: clustering is worth attempting when the joint
distribution of the variables is \emph{multipeaked}, and the first decision is
whether the data are multimodal at all. If a distribution is unimodal, every
algorithm will return groups anyway, and every one of them will be an
artefact.

\begin{marginnote}[]
\entry{Multimodality}{A distribution with more than one peak. Clustering is
only meaningful if the peaks are physical --- here, if stars of different
birth sites occupy different regions of C-space.}
\entry{No free lunch}{No algorithm is best on every problem: each encodes
assumptions about shape, scale, density and metric. Choosing a method is
choosing which assumptions are safe for your data.}
\end{marginnote}

\subsection{Distance in C-space}

Every algorithm in this workbook reduces to a similarity judgement, and all of
them use either Euclidean or cosine distance. Two transformations in
\textbf{Section~\ref{sec:data}} make those two compatible, and one
transformation is a genuine modelling choice:

\[
d_{\rm Euclid}(\mathbf{x},\mathbf{y}) = \left\| \mathbf{x} - \mathbf{y}
\right\|_2, \qquad
d_{\rm cos}(\mathbf{x},\mathbf{y}) = 1 - \frac{\mathbf{x}\cdot\mathbf{y}}
{\|\mathbf{x}\|_2\,\|\mathbf{y}\|_2}.
\]

After per-element standardisation and per-star $L_2$ normalisation, all stars
lie on the unit sphere in 16 dimensions and $d_{\rm Euclid}$ is a monotone
function of $d_{\rm cos}$ --- so t-SNE and UMAP (Euclidean internally) and
EVoC (cosine) are looking at the same geometry, and comparing them is fair.
The Euclidean alternative without normalisation is not wrong, but it inherits
a scale from the standardisation: elements with a long tail of outliers become
the elements that decide every distance.

High dimension does the rest. Two effects matter here.

\begin{enumerate}
\item \textbf{Concentration.} In $d$ dimensions, the ratio between the
  furthest and nearest neighbour distance shrinks as $d$ grows: distances
  \emph{concentrate} around a common value. A density threshold that worked in
  2-D has to be re-tuned, and the contrast between ``near'' and ``far''
  degrades. This is the everyday experience of the curse of dimensionality in
  clustering: the kNN graph of 16-D abundances has edges that are only
  slightly shorter than the edges you would get by chance.
\item \textbf{Dilution by concatenation.} Adding dimensions that carry no
  cluster information can only reduce the relative contribution of the
  dimensions that do. The consequence in this project is concrete. On the same
  982 member stars, in the same pipeline
  (\textbf{Section~\ref{sec:benchmark}}), the cluster-only homogeneity is
  $0.56$ from 16 abundances, $0.94$ from 4 kinematic coordinates, and
  $0.75$ from the 20-dimensional concatenation of the two. The blend is
  \emph{worse} than the kinematics alone: the chemistry did not add
  information, it added dimensions.
\end{enumerate}

That last result deserves emphasis because the intuitive fix --- ``if both
signals are useful, feed both'' --- is the one most people try first:

\begin{issues}[THE TRAP OF CONCATENATION]
Fusing two signals of very different dimensionality into one Euclidean space
is not a fusion of information; it is a weighted sum in which the weight of
each channel is set by how many coordinates it happens to contribute. A
256-dimensional spectral latent next to 4 kinematic coordinates gives the
spectrum 98\% of the vote, whatever the physics says. The two ways out are to
reduce the high-dimensional channel before concatenating it, or to combine the
two at the level of the \emph{decisions} rather than the features --- the
two-stage pipeline of \S\ref{sec:benchmark}, which uses kinematics for
recall and a spectral distance for precision.
\end{issues}

\subsection{Four tasks that are easily confused}

The literature --- and our own first drafts --- mixes up four different
questions that all get called ``clustering''. Keeping them separate is the
single most useful discipline in this project.

\begin{enumerate}
\item \textbf{Membership determination.} For a given cluster, label each star
  in a field \emph{member} or \emph{field}. This is the operationally useful
  task, and it is what a survey pipeline actually needs.
\item \textbf{Discovery.} Blind search: find groups without knowing in advance
  where to look or what they should look like. This is what chemical tagging
  means in the ``strong'' sense of \S\ref{sec:intro}, and it is much harder:
  the algorithm must invent the groups as well as find their members.
\item \textbf{Refinement.} Clean an existing candidate list --- remove field
  contaminants from a known cluster's sample. This is the task our two-stage
  pipeline performs, and it is measurably easier than discovery.
\item \textbf{Cluster-only separation.} Take the members of several clusters,
  with no field stars present, and ask whether the clusters are
  distinguishable from one another. This is the experiment of
  \citet{GarciaDias:19}, recreated in \S\ref{sec:benchmark} as the
  \emph{baseline}. Its score is an upper bound on what any field-level search
  could do: if clusters cannot be told apart when nothing else is present,
  they cannot be recovered from a field either.
\end{enumerate}

\begin{table}[htbp]
\tabcolsep6pt
\caption{The assumptions each method in this workbook makes explicit. Read it
as a checklist: when a method fails on our data, one of these entries is
usually the reason, and the next chapter takes over because of it.}
\label{tab:assumptions}
\begin{center}
\begin{tabular}{@{}p{4.6pc}p{8.2pc}p{9.6pc}@{}}
\hline
Method & Assumes & Fails when\\
\hline
K-means & clusters are convex, similar in size, separable by a centre;
  $K$ known & shapes are elongated or nested; outliers pull centres\\
$k$NN & a metric is meaningful and local density is informative &
  dimensionality is high; the $k$ is chosen without looking at the data\\
DBSCAN & one density threshold separates cluster from noise &
  clusters and noise share a density; densities vary between clusters\\
HDBSCAN$^*$ & density varies, but a single minimum cluster size applies &
  clusters are smaller than the floor; the metric is misleading\\
PLSCAN & persistence, not size, identifies a cluster &
  clusters are genuinely short-lived in scale-space (rare)\\
t-SNE & local neighbourhoods are what matter &
  inter-cluster distances are being interpreted\\
UMAP & the manifold is locally uniform &
  global geometry is treated as metric information\\
EVoC & a cosine-style graph and one internal embedding serve all scales &
  the metric is wrong for the data\\
\hline
\end{tabular}
\end{center}
\end{table}

\subsection{Validation, in outline}

Because there are no labels, validation is a design decision, not a
calculation. The toolbox has two halves, developed in
\textbf{Section~\ref{sec:validation}}:

\begin{itemize}
\item \textbf{Internal} indices --- sum of squared errors, silhouette, the dip
  test --- measure how \emph{tidy} a partition is. They are computed from the
  data alone, they are cheap, and they cannot tell you whether the partition
  is real.
\item \textbf{External} indices compare the partition with something the
  algorithm never saw. In this project that something is kinematic membership;
  in general it is any independent measurement.
\end{itemize}

The dangerous middle is a number that looks external but is not: a metric that
uses the quantity you are trying to predict, or a baseline that has collapsed
to two groups and scores well for the wrong reason. Both appear in our own
results and are flagged there.

\begin{issues}[EXERCISES]
\begin{enumerate}
\item The concentration of distances can be demonstrated in five lines:
  sample $n=1000$ points uniformly in $\mathbb{R}^d$, compute for each point
  the ratio of the distance to its nearest and to its furthest neighbour, and
  plot the mean ratio against $d$ for $d = 1,\dots,50$. At what $d$ does the
  nearest neighbour stop being meaningfully nearer than the furthest one?
\item Standardise the 16-D matrix, then row-normalise it. A star with an
  unusually low $[$Fe/H$]$ and a star with an unusually high one now have the
  same norm. What did normalisation do to the information about overall
  metallicity, and why is that mostly a feature here rather than a bug?
\item Pick three clusters from \textbf{Table~\ref{tab:clusters}} --- one with
  more than 60 members, one with about 20, one with fewer than 10 --- and
  write down what ``$K$ is known'' would even mean for a field of 25 clusters
  with these counts. Which of the four tasks of \S 3.3 does a fixed $K$
  violate?
\end{enumerate}
\end{issues}
